\chapter*{Tóm tắt}
\label{summary}

Trong bối cảnh học tập hiện đại, người học thường phải sử dụng đồng thời nhiều công cụ rời rạc để ghi chú, ôn tập, quản lý thời gian và tra cứu kiến thức, khiến dữ liệu học tập bị phân tán, khó liên kết và khó duy trì thói quen ôn tập chủ động. Các nền tảng hiện có phần lớn chỉ giải quyết tốt một nhu cầu riêng lẻ, chưa hỗ trợ toàn diện cả quá trình học từ tiếp nhận, tổ chức đến ôn tập và đánh giá kiến thức. Vì vậy, việc xây dựng một hệ thống học tập tích hợp ứng dụng trí tuệ nhân tạo là cần thiết nhằm hỗ trợ người học quản lý và nâng cao hiệu quả tự học một cách chủ động.

Mục tiêu của đề tài là phát triển một hệ thống học tập thông minh đa nền tảng, cho phép người dùng quản lý tài liệu, ghi chú, Flashcard, bài kiểm tra, mục tiêu và thời gian học tập trong cùng một môi trường thống nhất.

Đề tài kết hợp mô hình ngôn ngữ lớn (LLM - Large Language Model) với các agent chuyên biệt đảm nhiệm từng tác vụ học tập như tóm tắt tài liệu, sinh Flashcard, sinh bài kiểm tra và đề xuất lộ trình học tập cá nhân hóa. Việc ôn tập được cá nhân hóa theo thuật toán lặp cách quãng dựa trên mức độ ghi nhớ của từng người học, trong khi tính năng cộng tác ghi chú cho phép nhiều người cùng chỉnh sửa và đồng bộ nội dung theo thời gian thực. Hệ thống được đánh giá thông qua kiểm thử chức năng toàn diện và triển khai thử nghiệm thực tế với người dùng.

Chúng tôi đã xây dựng một hệ thống hỗ trợ học tập hoàn chỉnh, gồm 2 phân hệ : ứng dụng Web và ứng dụng Mobile, đã được triển khai và đưa vào sử dụng thực tế, hỗ trợ đầy đủ các chức năng ghi chú, Flashcard, bài kiểm tra, quản lý mục tiêu và thời gian học tập. Kết quả kiểm thử cho thấy hệ thống hoạt động ổn định và đáp ứng tốt các mục tiêu đề ra, đồng thời cho thấy tiềm năng ứng dụng thực tiễn của trí tuệ nhân tạo trong việc hỗ trợ học tập cá nhân hóa.
